\subsection{Augmented Assignment}

Augmented assignment combines an operator with assignment: \texttt{x += 1} instead of \texttt{x = x + 1}. The syntax suggests in-place mutation familiar from C, but for immutable singular types the reality is almost the opposite.

\subsubsection{Numeric Augmented Assignment (\texttt{+=}, \texttt{-=}, \texttt{*=}, \texttt{/=}) as Read-Operate-Rebind for Immutable Singular Values}

Compare C and Python increment at the machine level.

In C, \texttt{x++} or \texttt{x += 1} on a stack-allocated \texttt{int} typically compiles to load-add-store instructions that modify the existing memory word in place. No new integer object is allocated; no garbage is produced.

In Python, \texttt{x += 1} on an immutable \texttt{int} performs a three-stage heap operation:

\begin{enumerate}
    \item \textbf{Read:} dereference the name \texttt{x} to obtain the current \texttt{PyLongObject};
    \item \textbf{Operate:} dispatch \texttt{PyNumber\_Add} (or the \texttt{\_\_iadd\_\_}/\texttt{\_\_add\_\_} protocol), allocating a \emph{brand-new} \texttt{PyLongObject} for the result;
    \item \textbf{Rebind:} store the new object's pointer in the namespace slot for \texttt{x}, decrementing the reference count on the old object.
\end{enumerate}

The old integer object is not edited---integers are immutable. The name is rebound to a new heap object. The discarded object becomes garbage unless other references or small-integer caching keep it alive:

\begin{verbatim}
import sys

x = 42
old_id = id(x)
old_size = sys.getsizeof(x)

x += 1

# id(x) != old_id
# sys.getsizeof may differ for large magnitudes
# the PyLongObject formerly bound to x awaits refcount reclamation
\end{verbatim}

This is a deceptive illusion for C programmers: \texttt{+=} \emph{looks} like compound assignment into existing storage, but on immutable types it is syntactic compression for read-operate-rebind with mandatory heap allocation of the result object.

The same pattern applies to \texttt{-=}, \texttt{*=}, \texttt{/=}, and other augmented forms on \texttt{int}, \texttt{float}, \texttt{bool}, \texttt{str}, and other immutable singular domains. Division deserves special mention: \texttt{number /= 5} may rebind \texttt{number} from \texttt{int} to \texttt{float} because \texttt{number / 5} produces a \texttt{PyFloatObject}, not because the integer was mutated.

For immutable singular values:

\begin{quote}
Augmented assignment reads the current object, performs the operation (allocating a result object), and rebinds the name to that result. It does not modify bits in place.
\end{quote}

\subsubsection{Mutation vs. Rebinding Preview for Later Mutable Containers}

The preceding rule is reliable for immutable singular values. Mutable containers introduce a different code path: types such as \texttt{list} implement \texttt{\_\_iadd\_\_} to mutate the existing object in place when possible.

\begin{verbatim}
numbers = [10, 20]
alias = numbers
numbers += [30]       # mutates the shared list object in place
# alias also sees [10, 20, 30]; id(numbers) == id(alias)

numbers = numbers + [30]  # creates new list; rebinds numbers only
# alias still sees [10, 20]
\end{verbatim}

For numbers and strings, read augmented assignment as read-operate-rebind with heap allocation. For mutable containers, do not assume this; inspect whether the type implements in-place \texttt{\_\_iadd\_\_} or equivalent mutating behavior.
